\section{Limitations of this Survey}
\label{sec:limitations}

We state the boundaries of the survey explicitly. First, our organising axis, \emph{degree
of self-improvement}, is deliberately code-centric; it foregrounds systems that emit and
revise programs and treats weight-based policies (\S\ref{sec:vla}) and reinforcement-learning
skill discovery (\S\ref{sec:skill}) as context rather than as the object of the taxonomy.
A survey centred on representation learning or on real-world data collection would draw the
map differently. Second, the frontier we emphasise (\S3.1.5) is populated by very recent,
in some cases concurrent, systems; benchmark numbers across them are not directly
comparable, and we therefore report capabilities qualitatively rather than ranking methods.
Third, the field moves quickly: several 2025--2026 systems we cite were released while this
survey was being written, and the coverage should be read as a snapshot. Finally, the
``skill economy'' framing (\S\ref{sec:open}) is grounded in an emerging commercial trend
(one vendor's marketplace at the time of writing); we treat it as motivation for open
problems, not as a mature body of results.
